\section{实验设计与验证计划}

本节目前只给出实验框架和写作口径，不填写尚未获得的实验结果。实验章节的作用不是形成 planner 排行榜，而是围绕本文 claim 组织证据：\textbf{显式晃液状态预测是否能在单纯平滑控制之外进一步降低液体晃动，并且不以不可接受的路径跟踪、到达时间或实时性代价为条件。} 后续完成仿真或实物实验后，应在本节相应位置填入数据、图表和失败案例，而不是改变主线问题。

\subsection{实验问题与证据链}

根据 Introduction 和 Method 的论证，实验需要回答四类问题。\cref{tab:experiment_claim_map} 给出每类问题对应的证据形式。表中 ``待填'' 表示后续实验完成后补充具体数据、图号或表号。

\begin{table}[t]
  \centering
  \footnotesize
  \setlength{\tabcolsep}{3pt}
  \caption{实验问题、claim 与所需证据}
  \label{tab:experiment_claim_map}
  \begin{tabular}{p{0.28\linewidth} p{0.32\linewidth} p{0.28\linewidth}}
    \toprule
    实验问题 & 支持的 claim & 需要的证据 \\
    \midrule
    smooth-only 是否足够防晃 & 液体晃动不是单纯轨迹平滑问题 & \texttt{B\_ours} vs. \texttt{B\_smooth}，待填 \\
    晃液状态预测是否有独立贡献 & 低阶晃液状态提供动态记忆 & \texttt{B\_slosh} vs. \texttt{B0}，待填 \\
    完整方法是否保持 local planner 能力 & \SPMPC{} 同时优化路径进度、底盘控制和液体响应 & tracking、travel time、success，待填 \\
    普通 online planner 是否缺少液体状态 & ordinary local planner online but not liquid-aware & 外部 baseline 对比，待填 \\
    \bottomrule
  \end{tabular}
\end{table}

所有结果段落后续建议使用统一结构：
\begin{equation}
  \text{Claim} \rightarrow \text{Evidence} \rightarrow
  \text{Interpretation} \rightarrow \text{Limitation}.
\end{equation}
也就是说，每个图表都必须说明它支持哪个 claim、排除了哪个替代解释，以及仍然有哪些边界。

\subsection{实验场景与公平协议}

第一版实验范围应与本文问题定义一致：给定安全参考路径上的移动底盘开口液体运输。实验场景可以分为仿真验证和实物验证两层。仿真用于快速检查不同 planner 在同一参考路径和约束下的行为；实物实验用于评价真实液面响应，并验证内部 model proxy 是否具有解释价值。

公平协议应在所有正式对比前固定：

\begin{enumerate}
  \item \textbf{统一任务：} 相同起点、终点、参考路径、容器配置和任务终止条件；
  \item \textbf{统一底盘限制：} 尽量对齐最大线速度、角速度、线加速度和角加速度；
  \item \textbf{统一初始化：} 正式仿真比较应使用 fresh simulation；若只是当前仿真诊断，应明确不计入正式结果；
  \item \textbf{统一评价：} 外部 baseline 不使用液面状态作为控制输入；真实液面指标来自 RGB 或其它外部观测；
  \item \textbf{统一失败判据：} 记录 timeout、未到达、碰撞、明显偏离参考路径、求解失败等情况；
  \item \textbf{统一记录：} 保存轨迹、速度、控制输入、solver time、内部 slosh proxy 和外部液面观测。
\end{enumerate}

这里需要特别区分两类量：\texttt{/spmpc/slosh\_height} 只能作为模型内部预测或诊断 proxy；真实液体结论必须依赖 RGB max-LCR、液面高度测量或其它外部观测指标。

\begin{figure}[t]
  \centering
  \figplaceholder{实验平台 / 评价流程图：\\Scout 或仿真底盘、容器、参考路径、RGB 相机、\\日志与液面指标计算流程。}
  \caption{实验平台与外部液面评价流程占位。后续应插入仿真/实物场景、相机视角、液面 ROI 和日志记录流程。}
  \label{fig:experiment_setup_placeholder}
\end{figure}

\subsection{内部消融实验}

内部消融用于回答“\SPMPC{} 的每个模块是否必要”。所有变体应共享同一 alpha-state \MPCC{} 框架、同一参考路径和同一底盘约束，只改变晃液模型/代价和平滑权重。计划比较的变体见 \cref{tab:ablation_plan}。

\begin{table}[t]
  \centering
  \footnotesize
  \setlength{\tabcolsep}{3pt}
  \caption{内部消融实验计划}
  \label{tab:ablation_plan}
  \begin{tabular}{p{0.18\linewidth} p{0.20\linewidth} p{0.18\linewidth} p{0.32\linewidth}}
    \toprule
    变体 & 晃液模型/代价 & 平滑约束 & 主要回答的问题 \\
    \midrule
    \texttt{B0} & 无 & 弱 & 基础 \MPCC{} 的路径跟踪和液体响应基线 \\
    \texttt{B\_smooth} & 无 & 强 & 只增强 smoothness 是否足够防晃 \\
    \texttt{B\_slosh} & 有 & 弱 & 显式晃液状态预测是否有独立贡献 \\
    \texttt{B\_ours} & 有 & 强 & 完整方法是否兼顾防晃、跟踪和实时性 \\
    \bottomrule
  \end{tabular}
\end{table}

后续结果解释应重点围绕四组比较展开：\texttt{B\_smooth} vs. \texttt{B0} 用于量化平滑性收益；\texttt{B\_slosh} vs. \texttt{B0} 用于量化液体状态预测收益；\texttt{B\_ours} vs. \texttt{B\_smooth} 用于回答 smooth-only 是否足够；\texttt{B\_ours} vs. \texttt{B\_slosh} 用于说明平滑控制仍然是可执行局部规划的一部分。

\begin{figure}[t]
  \centering
  \figplaceholder{消融结果示例图：\\同一路径下 \texttt{B\_smooth} 与 \texttt{B\_ours} 的速度、\\预测/真实液面响应和路径跟踪曲线。}
  \caption{smooth-only 与 slosh-aware 消融对比图占位。后续应插入代表性试次曲线，用于展示相近平滑度或相近到达时间下的液面响应差异。}
  \label{fig:ablation_example_placeholder}
\end{figure}

\subsection{外部 local planner 对比}

外部对比用于回答“普通 online local planner 调好以后是否已经足够”。这一组 baseline 不是为了证明 \SPMPC{} 在所有导航指标上排名第一，而是为了提供 gap evidence：普通 planner 可以在线生成可执行和平滑运动，但通常不传播液体动态记忆。

候选 baseline 分为两类。第一类是同层 mobile-robot local planner，例如 DWA、TEB、\texttt{mpc\_local\_planner} 或 LT-DWA-family planner；它们是外部主对比对象。第二类是 anti-slosh profile 或 offline trajectory reference，例如 Hamaguchi-style velocity/path profile 或 Lim-style slosh-constrained trajectory；它们用于解释相关工作差异，不一定作为同层 online baseline 进入主表。最终进入主表的方法必须满足接口可运行、约束可对齐、任务定义一致和日志可复现四个条件。

\begin{table}[t]
  \centering
  \footnotesize
  \setlength{\tabcolsep}{3pt}
  \caption{外部对比方法的角色}
  \label{tab:external_baselines}
  \begin{tabular}{p{0.25\linewidth} p{0.30\linewidth} p{0.31\linewidth}}
    \toprule
    方法类型 & 示例 & 在论文中的作用 \\
    \midrule
    普通在线局部规划器 & DWA / TEB / non-slosh MPC / LT-DWA & 同层 baseline：online but not liquid-aware \\
    防晃 profile 或离线轨迹 & Hamaguchi-style / Lim-style reference & 相关工作参照：liquid-aware but not same local-planner layer \\
    \SPMPC{} 内部变体 & \texttt{B0}, \texttt{B\_smooth}, \texttt{B\_slosh}, \texttt{B\_ours} & 模块消融：分离 smooth 与 slosh-state contribution \\
    \bottomrule
  \end{tabular}
\end{table}

\subsection{评价指标}

评价指标应覆盖液体运动、任务完成、路径跟踪、控制平滑性和在线计算五个方面，如 \cref{tab:metrics_plan} 所示。填结果时应同时报告均值、方差或置信区间，并保留失败试次；如果样本数有限，则应明确写成 pilot evidence，而不是强结论。

\begin{table}[t]
  \centering
  \footnotesize
  \setlength{\tabcolsep}{3pt}
  \caption{评价指标与叙事作用}
  \label{tab:metrics_plan}
  \begin{tabular}{p{0.28\linewidth} p{0.34\linewidth} p{0.26\linewidth}}
    \toprule
    指标类别 & 具体指标 & 叙事作用 \\
    \midrule
    真实液面运动 & RGB max-LCR、RMS-LCR、残余振荡 & 验证是否真的降低晃动 \\
    模型诊断 & predicted slosh state、\texttt{/spmpc/slosh\_height} & 解释模型预测，不能替代真实评价 \\
    任务完成 & success、timeout、travel time & 排除“只是变慢/停下”的解释 \\
    路径跟踪 & contour error、path deviation & 验证仍是合格 local planner \\
    控制平滑 & velocity、acceleration、angular acceleration RMS & 区分 smoothness 与 slosh-aware 的作用 \\
    在线计算 & solver time、control frequency、failure count & 验证 receding-horizon 可实时运行 \\
    \bottomrule
  \end{tabular}
\end{table}

\subsection{结果小节占位结构}

实验完成后，本节可以按下面结构展开。当前版本只保留占位，避免在没有数据前写结果性判断。

\begin{enumerate}
  \item \textbf{Setup and fairness check：} 说明路径、容器、底盘限制、初始化、日志和评价方法是否一致；
  \item \textbf{Internal ablation：} 填入 \texttt{B0}/\texttt{B\_smooth}/\texttt{B\_slosh}/\texttt{B\_ours} 的主要表格和典型曲线；
  \item \textbf{External planner comparison：} 填入普通 local planner baseline 的对比表和代表性失败/成功案例；
  \item \textbf{Real-liquid observation：} 填入 RGB 或外部液面测量结果，并与内部 proxy 分开解释；
  \item \textbf{Real-time analysis：} 填入 solver time、控制频率和求解失败统计；
  \item \textbf{Limitations：} 说明模型误差、权重敏感性、路径类型、液体参数和传感噪声等边界。
\end{enumerate}

最终写作时，若某个实验没有完成，应明确标注为未完成或移到 future work，不能用内部 proxy 替代真实液面结论，也不能把 planned comparison 写成 observed result。